% ---------------------------------------------------------------------------
% Production / assembly / distribution network section (Pirhooshyaran & Snyder).
%
% INPUT-ABLE FRAGMENT: this file is \input{} into the manuscript
% learning_inventory_control_policies_es.tex and reuses ITS preamble macros and
% TikZ styles (envbox, envtitle, envsupply/envstock/envsink/envlost, reqflow/
% matflow/demandflow/lostflow, ltlabel/reqlabel, costH/costP, the boxed
% "Per-period event sequence" panel, booktabs, \ding). It defines NO new styles
% and contains NO \documentclass.
%
% BIBLIOGRAPHY: this section cites \citet{pirhooshyaran2021network}
% (Pirhooshyaran & Snyder 2021, arXiv:2006.05608), which is NOT yet in
% paper/references.bib; the entry must be added before compiling (out of scope
% of this fragment).
%
% HONESTY STATUS (load-bearing -- do not soften): this environment is FAITHFUL
% (verified equation-by-equation against the paper, eqs. 1-13) but
% literature_verified=FALSE. The ONLY published number it reproduces is the
% single-node newsvendor cost 127.11 (env exact DP -> ~127.10, <1%), which
% certifies env DYNAMICS, not the case3 result. The case3 comparator is the
% environment's OWN best grid-searched pairwise base-stock, a RESEARCH baseline,
% NOT a published optimum. The serial optimum 47.65 is STRUCTURALLY UNREACHABLE
% by this env (echelon-vs-local position mismatch) and lives in the serial
% section; it is NOT a target here. This is a learned-beats-own-best-heuristic
% research result, never a published-number beat.
%
% DATA-PROVENANCE NOTE:
%   Instance + dynamics : src/problems/multi_echelon/
%       production_assembly_distribution_network/ (env.rs: process-all-raw-on-
%       arrival, proportional allocation with backorder carry-over, order-after-
%       demand, holding on raw+finished+in-transit per cost eq.3),
%       heuristics/base_stock.rs (pairwise_base_stock_requests, eq.5),
%       literature/references.rs (PRIMARY_REFERENCE_INSTANCE
%       pirhooshyaran2021_serial_case3; literature_verified=false).
%   Results             : scripts/production_assembly_distribution_network/
%       RESULTS_case3_learned_vs_own_best_heuristic.md (committed 14b45d4, full
%       budget, 4000 held-out paired-CRN paths). Gate = env's own best pairwise
%       base-stock OUL [8,7,9], held-out 60.24/period. Learned depth-2 oblique
%       linear: seed123 57.25 +/- 0.22 (-4.96%), seed321 54.96 +/- 0.23
%       (-8.77%); depth-3 seed123 57.85 +/- 0.25 (-3.97%). Gen-0 flow warm-start
%       70.85 +/- 0.61.
% ---------------------------------------------------------------------------
\section{Production / assembly / distribution network}
\label{sec:pirhoo}

\subsection{Problem formulation}
\label{sec:pirhoo-formulation}

Figure~\ref{fig:env-pirhoo} depicts this environment as a flow network --- the per-relation order
requests, the lead-time shipments of raw material between nodes, the ``process all raw on arrival''
conversion to finished goods, the external demand at the downstream node, and where holding (on raw,
finished, and in-transit stock) and the backorder penalty are charged.

\begin{figure}[t]
  \centering
  \resizebox{\textwidth}{!}{%
  \begin{tikzpicture}[x=1cm,y=1cm]
    % Serial case-3 instantiation of the general supply-network MDP: a 3-node
    % serial chain 0 -> 1 -> 2 (node 0 sole source), external demand at node 2.
    % Each node carries raw + finished inventory; order requests UP (dashed
    % control) on each supply relation, lead-time material DOWN (solid). Backlog
    % is carried over at the demand node.
    \node[envsupply] (sup) at (0,0)      {external\\source};
    \node[envstock]  (n0)  at (5.4,0)    {node $0$\\raw\,$\mid$\,finished};
    \node[envstock]  (n1)  at (11.0,0)   {node $1$\\raw\,$\mid$\,finished};
    \node[envstock]  (n2)  at (16.6,0)   {node $2$\\raw\,$\mid$\,finished};
    \node[envsink]   (dem) at (22.0,0)   {demand\\$D_t\sim N(5,1)$};
    \node[envlost]   (bo)  at (22.0,-3.6){backorder\\(carried)};

    % per-relation requests UP (dashed control)
    \draw[reqflow] (n0.north west) to[bend right=26] (sup.north east);
    \draw[reqflow] (n1.north west) to[bend right=26] (n0.north east);
    \draw[reqflow] (n2.north west) to[bend right=26] (n1.north east);
    \node[reqlabel] at (10.9,1.95) {\ding{183}\, pairwise order-up-to on each relation};

    % lead-time material DOWN (solid)
    \draw[matflow] (sup.south east) to[bend right=10] (n0.west);
    \draw[matflow] (n0.south east)  to[bend right=10] (n1.west);
    \draw[matflow] (n1.south east)  to[bend right=10] (n2.west);
    \node[ltlabel] at (8.1,-1.5) {\ding{185}\, ship (lead $L$); \emph{process all raw on arrival}};

    \draw[demandflow] (n2) -- (dem)
          node[midway,above=3pt,ltlabel] {\ding{186}\, demand};
    \draw[lostflow] (dem) -- (bo);

    \node[costH,above=8mm of n1.north] {holding on raw\,$+$\,finished\,$+$\,in-transit};
    \node[costP,right=2pt of bo.east] {penalty $b\,B_t$};

    \begin{scope}[on background layer]
      \node[envbox,fit=(current bounding box)] (box) {};
    \end{scope}
    \node[envtitle,above=3pt of box.north,anchor=south]
          {General supply network (serial case~3 instantiation; external demand at the downstream node)};
  \end{tikzpicture}%
  }
  \par\vspace{6pt}\centerline{%
    \tikz\node[draw=mydarkblue!55,rounded corners=4pt,fill=mydarkblue!4,inner sep=7pt]{%
      \begin{minipage}{\dimexpr\linewidth-18pt\relax}\centering
        {\small\bfseries\color{mydarkblue!85}Per-period event sequence}\\[4pt]
        {\small\begin{tabular}{@{}p{4.3cm}p{4.3cm}p{4.3cm}@{}}
          \ding{182}\,observe per-node raw/finished inventory, pipelines, backlog & \ding{183}\,place a pairwise order-up-to request on each supply relation & \ding{184}\,external demand $D_t$ is observed at the downstream node\\[4pt]
          \ding{185}\,shipments placed $L$ ago arrive; each node \emph{processes all raw on arrival} into finished goods & \ding{186}\,downstream demand served from finished stock; shortfall \emph{backordered} and carried over & \ding{187}\,charge holding on raw, finished, and in-transit stock (cost eq.~3) and penalty $b$ on backorders\\
        \end{tabular}}
      \end{minipage}};}\par
  \caption{General production/assembly/distribution supply-network environment of
  \citet{pirhooshyaran2021network}, drawn as an event-based interaction; we instantiate it at the
  paper's serial case~3 (a 3-node serial chain $0\!\to\!1\!\to\!2$, node $0$ the sole source, external
  demand $N(5,1)$ at the downstream node $2$). Each node carries raw-material and finished-goods
  inventory. Each period a \emph{pairwise order-up-to} request is placed on every supply relation (thin
  dashed control arrows); shipments placed $L$ periods earlier arrive (solid material arrows) and each
  node \emph{processes all raw material on arrival} into finished goods. Downstream demand is served
  from finished stock and any shortfall is \emph{backordered and carried over}. Holding cost is charged
  on raw, finished, \emph{and} in-transit stock (the paper's cost equation~3) and a penalty $b$ on
  backorders; the boxed panel gives the per-period event sequence
  (Section~\ref{sec:pirhoo-formulation}, Eq.~\eqref{eq:pirhoo-cost}).}
  \label{fig:env-pirhoo}
\end{figure}

We consider the general production/assembly/distribution supply-network model of
\citet{pirhooshyaran2021network}: a directed network of nodes, each carrying raw-material and
finished-goods inventory, connected by supply relations with deterministic shipment lead times. The
control is a pairwise order-up-to policy on the supply relations; each node \emph{processes all raw
material on arrival} into finished goods, external demand arrives at the downstream node(s), and
shortfalls are backordered and carried over. We instantiate it at the paper's \emph{serial case~3}
(a 3-node serial chain), the instance wired into the environment. Our implementation is verified
\emph{equation-by-equation} against the paper (cost equation~3; a worked-transition fixture reproduces
the per-period cost by hand), but it is \emph{not} literature-verified at the network level: see the
honesty note below.

\paragraph{Timing of a period.}
At the start of period $t$ the controller observes each node's raw and finished inventory, the
in-transit pipelines, and the downstream backlog, and places a pairwise order-up-to request on each
supply relation. External demand is observed; shipments placed $L$ periods earlier arrive and each
node processes all arriving raw material into finished goods; downstream demand is served from finished
stock with proportional allocation and any shortfall is backordered; holding and penalty costs are
charged.

\paragraph{State.}
The state collects, for each node, its raw and finished inventory and in-transit pipeline, together
with the downstream backlog,
\begin{equation*}
    \mathbf{S}_t=\bigl(\{\text{raw}^{n}_t,\ \text{fin}^{n}_t,\ \text{pipeline}^{n}_t\}_{n},\ \text{backlog}_t\bigr).
\end{equation*}

\paragraph{Action.}
The decision is a per-relation order quantity (one entry per supply relation),
\begin{equation*}
    a_t=\bigl(q^{1}_t,\dots,q^{R}_t\bigr),
\end{equation*}
the order placed on each of the $R$ supply relations; for serial case~3, $R=3$ (the two internal edges
and the external source).

\paragraph{Transition.}
Orders placed on each relation are shipped subject to upstream availability and arrive after the
relation's lead time. On arrival each node processes all raw material into finished goods. Downstream
demand is served from finished stock; when finished stock is scarce it is allocated proportionally and
the unmet quantity is backordered and carried into the next period.

\paragraph{One-period cost.}
The cost incurred in period $t$ is the holding cost on \emph{all} inventory (raw, finished, and
in-transit) plus the backorder penalty at the demand node:
\begin{equation}
    \label{eq:pirhoo-cost}
    c_t \;=\; \sum_{n}\Bigl(h^{n}\bigl[\text{raw}^{n}_t+\text{fin}^{n}_t+\text{in-transit}^{n}_t\bigr]\Bigr) \;+\; b\,B_t ,
\end{equation}
with per-node holding rates $h^{n}$ and per-unit backorder penalty $b$ at the demand node ($B_t$ the
backorder level), following the paper's cost equation~3.

\paragraph{Objective.}
A stationary policy $\pi_{\btheta}$ maps the state to the per-relation orders. Performance is the
expected undiscounted average per-period cost over the horizon, $\min_{\btheta} C(\btheta)$, estimated
by simulation over shared demand-path seeds.

\FloatBarrier

\subsection{Policy}
\label{sec:pirhoo-policy}

This environment tests the recipe on a \emph{network} action space. \textbf{Input.} The policy input
is the per-node feature vector (raw and finished inventory, in-transit pipeline, and backlog),
divide-by-scale normalized. \textbf{Output.} A valid action is the per-relation order vector (a
\texttt{vector\_quantity} action of dimension $R$), clipped to the physical action box.
\textbf{Internals.} $\pi_{\btheta}$ uses the soft-tree backbone of Section~\ref{sec:soft-tree}: a
depth-2 tree with \emph{oblique} splits and \emph{linear} leaves, so each leaf maps the policy-state
features to the per-relation orders and oblique splits switch regime by inventory state. The design
choice we test is the \emph{leaf class}: a constant leaf can emit only a fixed order rate (and stays
at the flow regime), whereas a linear leaf reads finished inventory, in-transit pipeline, and backlog
to express inventory-position feedback on the same supply relations as the pairwise base-stock policy.
We warm-start CMA-ES at the steady-state flow rate (order the demand mean on each relation), the
simplest reproducible known-good point, and the search refines outward.

\subsection{Experiment setup}
\label{sec:pirhoo-setup}

We study the paper's serial case~3 (\texttt{pirhooshyaran2021\_serial\_case3}): a 3-node serial chain
$0\!\to\!1\!\to\!2$ with node $0$ the sole source, external demand $N(5,1)$ at node $2$ only, horizon
$T=10$, shipment lead times (external$\to$0, $0\to1$, $1\to2$) $=(2,1,1)$, local holding $[2,4,7]$, and
backorder cost $37.12$ at node $2$. The objective is the undiscounted average per-period cost.

\paragraph{Honesty status: a faithful but not literature-verified environment.}
This environment reproduces exactly one \emph{published} quantity: the single-node newsvendor cost of
the paper's Table~1 ($\mu=100,\sigma=10,h=10,p=30,L=1$), for which the environment's exact dynamic
program returns $\approx127.10$ against the published $127.11$ ($<1\%$). This certifies the
\emph{dynamics}. There is \emph{no} published optimum for the multi-node version of \emph{this} MDP, so
every multi-node reference instance carries \texttt{literature\_verified=false}. In particular, the
textbook serial optimum $47.65$ (Section~\ref{sec:serial}) is an \emph{echelon} base-stock cost and is
\emph{structurally unreachable} by this environment's \emph{local} pairwise policy (which targets the
raw-material position and excludes finished goods); we therefore do \emph{not} use $47.65$ (or any
published number) as the comparator here. The comparator is instead the environment's \emph{own} best
heuristic.

\paragraph{Comparator gate (a research baseline, not an optimum).}
The strongest heuristic native to this environment is the pairwise base-stock policy (the paper's
equation~5). We grid-search its per-relation order-up-to levels on a search-seed block and re-score the
argmin on a disjoint held-out block: the argmin is order-up-to $[8,7,9]$ with held-out cost
$60.24$ per period. We state explicitly that $60.24$ is the environment's own best heuristic, not an
optimum.

\paragraph{Evaluation protocol.}
Policies are trained with warm-started CMA-ES (population $24$, $60$ generations) and evaluated under
the common-random-number protocol of Section~\ref{sec:methods-eval}: learned policy and gate are scored
on the \emph{same} held-out block ($4000$ paths), paired on identical demand realizations.

\subsection{Results}
\label{sec:pirhoo-results}

Table~\ref{tab:pirhoo-results} reports the learned soft tree against the environment's own best
pairwise base-stock on the held-out block. The learned policy \emph{beats the environment's own best
heuristic}, robustly and reproducibly: the depth-2 oblique-linear tree improves the per-period cost
from $60.24$ to $57.25\pm0.22$ ($-4.96\%$) on one CMA seed and to $54.96\pm0.23$ ($-8.77\%$) on
another, and a depth-3 tree gives $57.85\pm0.25$ ($-3.97\%$); every configuration beats the gate by
more than $3.9\%$, outside the held-out standard error by $\ge 9\times$ SEM, across two seeds and two
depths. The flow warm-start (generation~0) sits at $70.85\pm0.61$, so the improvement is the CMA-ES
delta, not the initialization. The mechanism is a strictly richer control class on the \emph{same}
action relations: the pairwise base-stock policy reacts only to the local raw-material position, while
the learned linear-leaf tree additionally reads finished inventory, backlog, and inbound pipeline, and
switches behaviour by inventory regime; constant-leaf trees stay at the flow regime and lose to the
gate, confirming the leaf class is the lever.

\paragraph{Honest scope.}
This is a \emph{research} result on a faithful but \emph{not} literature-verified environment: the
learned policy beats the environment's \emph{own} best heuristic, not any published cost. We report it
as evidence that action design (the leaf class), not optimizer capacity, governs whether a black-box
search recovers structured-control performance --- consistent with the same finding on the lost-sales,
one-warehouse multi-retailer, and multi-echelon environments elsewhere in this work. Upgrading the
environment to literature-verified would require recovering the paper's exact order-up-to to
inventory-position simulation convention and binding the in-crate echelon-position policy; we leave
that to future work and do not present this row as a published-number benchmark.

\begin{table}[!htbp]
\centering
\small
\caption{Production/assembly/distribution network (serial case~3): best learned soft tree vs.\ the
environment's \emph{own} best pairwise base-stock gate, common-random-number evaluation ($4000$
held-out paths). Costs are undiscounted per-period averages (lower is better). $\Delta\%$ is the
learned policy's improvement over the gate; ``SEM mult.'' is that gain divided by the held-out
standard error. The gate ($60.24$) is the environment's own grid-searched heuristic, \emph{not} a
published optimum --- this is a research comparison, not a literature-number beat. Source: the
held-out ledger under \texttt{scripts/production\_assembly\_distribution\_network/}.}
\label{tab:pirhoo-results}
\begin{tabular}{lccc}
\toprule
Policy / config & Held-out cost ($\pm$ SEM) & $\Delta\%$ vs gate & SEM mult. \\
\midrule
Best pairwise base-stock (gate, OUL $[8,7,9]$) & $60.24$ & --- & --- \\
Flow warm-start (generation $0$)               & $70.85\pm0.61$ & $+17.6\%$ & --- \\
Learned soft tree, depth $2$, seed $123$       & $\mathbf{57.25\pm0.22}$ & $\mathbf{-4.96\%}$ & $\approx 14\times$ \\
Learned soft tree, depth $2$, seed $321$       & $\mathbf{54.96\pm0.23}$ & $\mathbf{-8.77\%}$ & $\approx 23\times$ \\
Learned soft tree, depth $3$, seed $123$       & $\mathbf{57.85\pm0.25}$ & $\mathbf{-3.97\%}$ & $\approx 10\times$ \\
\bottomrule
\end{tabular}
\par\medskip\footnotesize The comparator is the environment's own best heuristic on a faithful but
\texttt{literature\_verified=false} environment; the result is a research learned-beats-own-heuristic
comparison, not a published-number beat. The serial optimum $47.65$ (Section~\ref{sec:serial}) is
structurally unreachable by this environment's local pairwise policy and is not used as a target here.
\end{table}
\FloatBarrier
